\usepackage[a4paper,inner=24mm,outer=21mm,top=22mm,bottom=24mm,headheight=21pt]{geometry}
\usepackage[spanish,es-nodecimaldot]{babel}
\usepackage{fontspec}
\setmainfont{Libertinus Serif}
\setsansfont{Libertinus Sans}
\usepackage{unicode-math}
\setmathfont{Libertinus Math}
\usepackage{microtype}
\usepackage{xcolor}
\usepackage{graphicx}
\usepackage{booktabs}
\usepackage{tabularx}
\usepackage{longtable}
\usepackage{array}
\usepackage{enumitem}
\usepackage{needspace}
\usepackage{fancyhdr}
\usepackage{caption}
\usepackage{accsupp}
\usepackage{hyperref}
\usepackage{hyperxmp}
\usepackage{bookmark}
\usepackage[most]{tcolorbox}
\usepackage{fvextra}
\usepackage{tikz}
\usetikzlibrary{calc}

\definecolor{CourseNavy}{HTML}{24364B}
\definecolor{CourseBlue}{HTML}{446A8C}
\definecolor{CourseTeal}{HTML}{4F7F7B}
\definecolor{CourseAmber}{HTML}{A9713A}
\definecolor{CourseRed}{HTML}{965B59}
\definecolor{CourseInk}{HTML}{252A30}
\definecolor{CourseMuted}{HTML}{5D6670}
\definecolor{CourseRule}{HTML}{D8E0E7}
\definecolor{CoursePaper}{HTML}{FBFCFD}
\definecolor{CoursePaleBlue}{HTML}{F4F7FA}
\definecolor{CoursePaleTeal}{HTML}{F3F8F7}
\definecolor{CoursePaleAmber}{HTML}{FBF7F0}
\definecolor{CoursePaleRed}{HTML}{FAF4F3}
\definecolor{CourseCoverMist}{HTML}{EDF2F5}



\providecommand{\PSevenPdfTitle}{Cálculo Diferencial en una Variable}
\providecommand{\PSevenPdfSubject}{Problemas, prácticas y estudio}
\providecommand{\PSevenPdfIdentifierSuffix}{}
\providecommand{\PSevenPdfKeywords}{cálculo diferencial, ejercicios, prácticas, estudio}
\providecommand{\PSevenRunningTitle}{Problemas, prácticas y estudio}
\providecommand{\PSevenEdition}{Edición pública para estudiantes · v1.0}
\providecommand{\PSevenAccent}{CourseTeal}
\providecommand{\PSevenPale}{CoursePaleTeal}

\hypersetup{
  colorlinks=true,
  hypertexnames=false,
  linkcolor=CourseBlue,
  urlcolor=CourseTeal,
  pdftitle={\PSevenPdfTitle},
  pdfauthor={Giovanni Dalmasso, PhD},
  pdfsubject={\PSevenPdfSubject\PSevenPdfIdentifierSuffix},
  pdfkeywords={\PSevenPdfKeywords},
  pdfcreator={LuaLaTeX; fuente pública para estudiantes},
  pdfcopyright={Copyright © 2026 Giovanni Dalmasso},
  pdflicenseurl={https://creativecommons.org/licenses/by-sa/4.0/},
  pdfmetalang={es},
  pdflang={es-ES},
  pdfdisplaydoctitle=true,
  keeppdfinfo=true,
  bookmarksopen=true,
  bookmarksnumbered=true
}
\pdfstringdefDisableCommands{%
  \def\texttt#1{#1}\def\ensuremath#1{#1}%
  \def\PSevenTarget#1{}\def\PSevenID#1{}%
}

\setlength{\parindent}{0pt}
\setlength{\parskip}{4pt plus 1pt minus 1pt}
\setlist[itemize]{leftmargin=*,itemsep=2pt,topsep=3pt}
\setlist[enumerate]{leftmargin=*,itemsep=3pt,topsep=4pt}
\raggedbottom
\emergencystretch=3em
\tolerance=1500

\pagestyle{fancy}
\fancyhf{}
\fancyhead[L]{\small\sffamily\PSevenEdition}
\fancyhead[R]{\small\sffamily\PSevenRunningTitle}
\fancyfoot[C]{\thepage}
\renewcommand{\headrulewidth}{0.35pt}
\renewcommand{\headrule}{\hbox to\headwidth{\color{CourseRule}\leaders\hrule height \headrulewidth\hfill}}
\renewcommand{\chaptername}{Capítulo}

\newtcolorbox{PSevenNotice}[2][]{
  enhanced,breakable,frame hidden,boxrule=0pt,
  borderline west={1.8pt}{0pt}{\PSevenAccent},
  colback=\PSevenPale,coltitle=CourseInk,title={#2},fonttitle=\sffamily\bfseries,
  left=8pt,right=8pt,top=4pt,bottom=5pt,before skip=6pt,after skip=6pt,#1
}
\newtcolorbox{PSevenExercise}[2][]{
  enhanced,breakable,colback=white,colframe=CourseRule,boxrule=0.45pt,
  borderline west={2pt}{0pt}{\PSevenAccent},
  title={#2},fonttitle=\sffamily\bfseries,coltitle=CourseInk,
  left=8pt,right=8pt,top=5pt,bottom=6pt,before skip=7pt,after skip=7pt,#1
}
\newtcolorbox{PSevenAnswer}[2][]{
  enhanced,breakable,frame hidden,boxrule=0pt,
  borderline west={1.6pt}{0pt}{CourseTeal},colback=CoursePaleTeal,
  title={#2},fonttitle=\sffamily\bfseries,coltitle=CourseInk,
  left=8pt,right=8pt,top=4pt,bottom=5pt,before skip=6pt,after skip=6pt,#1
}
\newtcolorbox{PSevenSolution}[2][]{
  enhanced,breakable,frame hidden,boxrule=0pt,
  borderline west={1.8pt}{0pt}{CourseBlue},colback=CoursePaleBlue,
  title={#2},fonttitle=\sffamily\bfseries,coltitle=CourseInk,
  left=8pt,right=8pt,top=4pt,bottom=6pt,before skip=6pt,after skip=7pt,#1
}
\newtcolorbox{PSevenOptional}[2][]{
  enhanced,breakable,frame hidden,boxrule=0pt,
  borderline west={1.6pt}{0pt}{CourseAmber},colback=CoursePaleAmber,
  title={#2},fonttitle=\sffamily\bfseries,coltitle=CourseInk,
  left=8pt,right=8pt,top=4pt,bottom=5pt,before skip=6pt,after skip=6pt,#1
}
\newtcolorbox{PSevenExerciseFixed}[2][]{
  enhanced,colback=white,colframe=CourseRule,boxrule=0.45pt,
  borderline west={2pt}{0pt}{\PSevenAccent},
  title={#2},fonttitle=\sffamily\bfseries,coltitle=CourseInk,
  left=8pt,right=8pt,top=5pt,bottom=6pt,before skip=7pt,after skip=7pt,#1
}
\newtcolorbox{PSevenAnswerFixed}[2][]{
  enhanced,frame hidden,boxrule=0pt,
  borderline west={1.6pt}{0pt}{CourseTeal},colback=CoursePaleTeal,
  title={#2},fonttitle=\sffamily\bfseries,coltitle=CourseInk,
  left=8pt,right=8pt,top=4pt,bottom=5pt,before skip=6pt,after skip=6pt,#1
}
\newtcolorbox{PSevenSolutionFixed}[2][]{
  enhanced,frame hidden,boxrule=0pt,
  borderline west={1.8pt}{0pt}{CourseBlue},colback=CoursePaleBlue,
  title={#2},fonttitle=\sffamily\bfseries,coltitle=CourseInk,
  left=8pt,right=8pt,top=4pt,bottom=6pt,before skip=6pt,after skip=7pt,#1
}
\newtcolorbox{PSevenOptionalFixed}[2][]{
  enhanced,frame hidden,boxrule=0pt,
  borderline west={1.6pt}{0pt}{CourseAmber},colback=CoursePaleAmber,
  title={#2},fonttitle=\sffamily\bfseries,coltitle=CourseInk,
  left=8pt,right=8pt,top=4pt,bottom=5pt,before skip=6pt,after skip=6pt,#1
}

\newcommand{\PSevenMeta}[1]{\par\smallskip{\small\sffamily\color{CourseMuted}#1}\par\smallskip}
% Canonical IDs remain only in source anchors for traceability and are never typeset.
\newcommand{\PSevenID}[1]{}
\newcommand{\PSevenLink}[2]{\hyperlink{#1}{\textcolor{CourseBlue}{#2}}}
\makeatletter
\newcommand{\PSevenTarget}[1]{\Hy@raisedlink{\hyper@anchorstart{#1}\hyper@anchorend}}
\makeatother

\newcommand{\PSevenCoverArtPath}{src/problems/assets/cover.png}
\newcommand{\PSevenCoverArtwork}{%
  \BeginAccSupp{method=pdfstringdef,unicode,ActualText={Ilustración de cálculo diferencial con curvas, superficies, derivadas, circunferencia trigonométrica, datos y material de estudio}}%
  \includegraphics[height=0.47\textheight,keepaspectratio]{\PSevenCoverArtPath}%
  \EndAccSupp{}%
}

\newcommand{\PSevenCover}[3]{%
\begin{titlepage}
  \thispagestyle{empty}
  {\sffamily\color{CourseNavy}\fontsize{12}{15}\selectfont GIOVANNI DALMASSO}\par
  \vspace{0.55cm}
  {\color{CourseInk}\fontsize{29}{33}\selectfont\bfseries Cálculo Diferencial\\en una Variable\par}
  \vspace{0.28cm}
  {\sffamily\color{\PSevenAccent}\fontsize{18}{22}\selectfont\bfseries #1\par}
  \vspace{0.12cm}
  {\sffamily\color{CourseMuted}\large #2\par}
  \vspace{0.45cm}
  \centering\PSevenCoverArtwork\par
  \vspace{0.34cm}
  {\sffamily\bfseries\color{\PSevenAccent}#3\par}
  \vfill
  {\sffamily\small Versión 1.0\par}
\end{titlepage}
}

\newcommand{\PSevenTitleAndRights}[2]{%
\begin{titlepage}
  \thispagestyle{empty}
  \centering
  \vspace*{1.8cm}
  {\Huge\bfseries Cálculo Diferencial en una Variable\par}
  \vspace{0.45cm}
  {\Large\color{\PSevenAccent}#1\par}
  \vspace{0.45cm}
  {\large\sffamily #2\par}
  \vspace{1.05cm}
  {\Large \PublicAuthorDisplayName\par}
  \vspace{0.3cm}
  {\normalsize \PublicAuthorAffiliationLines\par}
  \vfill
  {\normalsize Versión 1.0 · 2026\par}
\end{titlepage}

\clearpage
\thispagestyle{empty}
\vspace*{0.24\textheight}
{\small
\textbf{Cálculo Diferencial en una Variable}\par
#1\par
#2\par
\PublicAuthorDisplayName\par
Versión 1.0, 2026\par
Handle DAU: \url{https://hdl.handle.net/20.500.14342/7179}\par
\vspace{1em}
© 2026 Giovanni Dalmasso. Salvo indicación en contrario, esta obra se
distribuye bajo la licencia Creative Commons Atribución-CompartirIgual 4.0
Internacional (CC BY-SA 4.0). Resumen legible:
\url{https://creativecommons.org/licenses/by-sa/4.0/}. Texto legal oficial en
español: \url{https://creativecommons.org/licenses/by-sa/4.0/legalcode.es}.\par
\vspace{0.7em}
Fotografía del autor: facilitada por la fotógrafa de IQS y utilizada con
permiso. Esta fotografía no está incluida en la licencia CC BY-SA 4.0 del
resto de la obra.\par
\vspace{0.7em}
Afiliación del autor, a efectos informativos: \PublicAuthorAffiliationInline.
Esta es una publicación del autor y no se presenta como publicación
institucional oficial.\par
\vspace{0.7em}
\textbf{Cita recomendada:} Dalmasso, G. (2026). \emph{Cálculo Diferencial en
una Variable — Problemas, prácticas y estudio} (versión 1.0).
\url{https://hdl.handle.net/20.500.14342/7179}\par
\vspace{0.7em}
\textbf{Ilustración de portada:} imagen generada con ChatGPT/OpenAI en 2026,
dirigida y seleccionada por Giovanni Dalmasso. El proyecto no conserva datos
verificables sobre la fecha exacta, el modelo, el prompt o la sesión de
generación.\par
}
\clearpage
}

\newcommand{\PSevenCompactTitle}[2]{%
  \thispagestyle{plain}
  {\sffamily\small\color{CourseMuted}Cálculo Diferencial en una Variable}\par
  {\LARGE\bfseries\color{CourseInk}#1}\par
  {\large\sffamily\color{\PSevenAccent}#2}\par
  \vspace{0.35cm}\color{CourseRule}\hrule\color{CourseInk}\vspace{0.25cm}
}

\newcommand{\PSevenLabFigure}[3]{%
  \par\medskip
  \begin{minipage}{\linewidth}
  \centering
  \BeginAccSupp{method=pdfstringdef,unicode,ActualText={#3}}%
  \includegraphics[width=0.88\linewidth]{#1}%
  \EndAccSupp{}%
  \captionof{figure}{#2}
  \end{minipage}
  \par\medskip
}
